\documentclass{article}
\date{}
\usepackage{hyperref}
\usepackage{../../apxproof}
\usepackage{amsthm}

\newtheoremrep{theorem}{Theorem}
\newtheoremrep*{lemma}[theorem]{Lemma}

\renewcommand{\appendixproofname}[2]{Proof of #1~#2}

\begin{document}
\section{Repeated and proof-deferred theorems}

\begin{theoremrep}\label{thm:first}
  Repeated theorem: this one is restated in the appendix.
\end{theoremrep}
\begin{proof}
  Proof of Theorem~\ref{thm:first}; appears in the appendix below the
  restated statement, titled by \texttt{\textbackslash appendixproofname}.
\end{proof}

\begin{lemmarep}\label{lem:second}
  Proof-deferred lemma: this one is \emph{not} restated in the appendix,
  but shares the theorem counter.
\end{lemmarep}
\begin{proof}
  Proof of Lemma~\ref{lem:second}; appears in the appendix without
  restating the lemma.
\end{proof}

\begin{theoremrep}\label{thm:third}
  A third theorem, restated.
\end{theoremrep}
\begin{proof}[An explicit title]
  Proof with an explicit optional title, which must override
  \texttt{\textbackslash appendixproofname}.
\end{proof}

\begin{theorem}\label{thm:plain}
  A plain (non-rep) theorem: stays in the main text only.
\end{theorem}
\begin{proof}
  Proof in the main text, not deferred.
\end{proof}

\section{Theorem without a proof}

\begin{lemmarep}\label{lem:noproof}
  A proof-deferred lemma with no proof at all; \texttt{apxproof} should
  not produce a proof for it in the appendix.
\end{lemmarep}
\noproofinappendix

\begin{theorem}
  Plain theorem after the orphan lemma.
\end{theorem}
\begin{proof}
  Proof in the main text.
\end{proof}

\end{document}
